\chapter{PENDAHULUAN}

\section{Latar Belakang}

% TODO: Tulis latar belakang penelitian Anda di sini.
% Jelaskan konteks, urgensi, dan relevansi topik yang dipilih.
% Rujuk pada penelitian terdahulu dan data sekunder jika ada.

Lorem ipsum dolor sit amet, consectetur adipiscing elit. Sed do eiusmod tempor incididunt ut labore et dolore magna aliqua. Ut enim ad minim veniam, quis nostrud exercitation ullamco laboris nisi ut aliquip ex ea commodo consequat. Duis aute irure dolor in reprehenderit in voluptate velit esse cillum dolore eu fugiat nulla pariatur. Excepteur sint occaecat cupidatat non proident, sunt in culpa qui officia deserunt mollit anim id est laborum.

Sed ut perspiciatis unde omnis iste natus error sit voluptatem accusantium doloremque laudantium, totam rem aperiam, eaque ipsa quae ab illo inventore veritatis et quasi architecto beatae vitae dicta sunt explicabo. Nemo enim ipsam voluptatem quia voluptas sit aspernatur aut odit aut fugit, sed quia consequuntur magni dolores eos qui ratione voluptatem sequi nesciunt.

\section{Rumusan Masalah}

% TODO: Tulis rumusan masalah dalam bentuk pertanyaan penelitian.
% Rumusan masalah harus spesifik, terukur, dan dapat dijawab.

Berdasarkan latar belakang di atas, rumusan masalah dalam penelitian ini adalah:

\begin{enumerate}
  \item Bagaimana ...?
  \item Bagaimana ...?
  \item Bagaimana ...?
\end{enumerate}

\section{Batasan Masalah}

% TODO: Tulis batasan masalah untuk memfokuskan penelitian.
% Batasan dapat mencakup ruang lingkup, teknologi, atau aspek yang tidak diteliti.

Agar penelitian ini tetap fokus dan terukur, batasan masalah yang ditetapkan adalah:

\begin{enumerate}
  \item Lorem ipsum dolor sit amet, consectetur adipiscing elit.
  \item Sed do eiusmod tempor incididunt ut labore et dolore magna aliqua.
  \item Ut enim ad minim veniam, quis nostrud exercitation ullamco laboris.
\end{enumerate}

\section{Tujuan Penelitian}

% TODO: Tulis tujuan penelitian yang ingin dicapai.
% Tujuan harus sesuai dengan rumusan masalah.

Tujuan dari penelitian ini adalah:

\begin{enumerate}
  \item Lorem ipsum dolor sit amet, consectetur adipiscing elit.
  \item Sed do eiusmod tempor incididunt ut labore et dolore magna aliqua.
  \item Ut enim ad minim veniam, quis nostrud exercitation ullamco laboris.
\end{enumerate}

\section{Manfaat Penelitian}

% TODO: Tulis manfaat penelitian bagi akademik dan praktis.

Penelitian ini diharapkan memberikan manfaat berikut:

\begin{enumerate}
  \item \textbf{Teoritis:} Lorem ipsum dolor sit amet, consectetur adipiscing elit.
  \item \textbf{Praktis:} Sed do eiusmod tempor incididunt ut labore et dolore magna aliqua.
\end{enumerate}

\section{Sistematika Penulisan}

% TODO: Sesuaikan sistematika penulisan dengan struktur penelitian Anda.

Sistematika penulisan proposal tugas akhir ini adalah sebagai berikut:

\begin{itemize}
  \item \textbf{Bab I:} Pendahuluan --- berisi latar belakang, rumusan masalah, batasan masalah, tujuan, manfaat, dan sistematika penulisan.
  \item \textbf{Bab II:} Penelitian Terkait --- membahas state of the art dan hasil-hasil penelitian sebelumnya yang relevan.
  \item \textbf{Bab III:} Metode dan Rancangan Penelitian --- menjelaskan metodologi, perancangan sistem, bagan alir, dan tahapan implementasi.
  \item \textbf{Bab IV:} Jadwal Penelitian --- memuat Gantt Chart berbasis bulan untuk perencanaan waktu penelitian.
\end{itemize}
